\workbookchapter{长上下文建模}

\begin{exercises}
\begin{exercise} 将 $T_0=4096$ 扩展为 $T_1=32768$，分别计算采用长度比与端点比的位置映射，比较最后一个位置的映射结果。说明为什么配置必须统一这两种约定。

\begin{solution}
长度比 $s_L=\frac{32768}{4096}=8$，末位置32767映为4095.875；端点比 $s_E=\frac{32767}{4095}$，映射 $f(i)=\frac{4095i}{32767}$，末位置恰为4095。两者差0.875个参考位置，均可作为明确约定，但不能在缓存、新查询或不同层间混用，否则相对相位不一致。
\end{solution}
\end{exercise}

\begin{exercise}\optional 对给定频率 $\omega=\frac{1}{256}$ 和距离 $\Delta=2048$，计算扩展因子为四时的插值相位，并写出固定查询与键条件下的点积分量变化上界。

\begin{solution}
原相位 $\Delta\omega=\frac{2048}{256}=8$，因子四插值后为2。固定q、k时精确旋转范数界为 $2\|\vect{q}\|\|\vect{k}\||\sin(\frac{2-8}{2})|=2|\sin3|\|\vect{q}\|\|\vect{k}\|$，还可用较松的 $6\|\vect{q}\|\|\vect{k}\|$。该界仅约束当前二维点积分量；不能直接约束深层网络，因为q、k也可能随扩展改变。
\end{solution}
\end{exercise}

\begin{exercise}\optional 从式\tbobject{eq:lc-base}推导频率比例，验证最高与最低频带的边界。解释为何 $d_r=2$ 不适用该基数规则。

\begin{solution}
若 $b'=b\,s^{\frac{d_r}{d_r-2}}$、$\omega_k=b^{-\frac{2k}{d_r}}$，则 $\frac{\omega'_k}{\omega_k}=s^{-\frac{2k}{d_r-2}}$。最高频$k=0$比例1，最低频 $k=\frac{d_r}{2}-1$ 比例 $\frac{1}{s}$。$d_r=2$ 时分母为零，且仅一个频带，无法同时满足“最高频不变”和“最低频缩为1/s”的非平凡要求。
\end{solution}
\end{exercise}

\begin{exercise}\optional 对式\tbobject{eq:lc-dilution}，给定目标权重 $0.9$，推导序列长度由101增至901时相关分数所需增加的量。说明真实多头网络与该构造相比多出了哪些变量。

\begin{solution}
设目标分数a、其余分数c，权重 $\frac{1}{[1+(T-1)\conste{}^{c-a}]}=.9$，解得 $a-c=\log[9(T-1)]$。T由101至901需要增加 $\log(\frac{900}{100})=\log9\approx2.197225$。真实模型的干扰分数不相等，多头值内容、层间表示和位置编码也变化；本构造揭示分母稀释，不是通用长度性能公式。
\end{solution}
\end{exercise}

\begin{exercise} 因果窗口包含最近八个位置，网络有六层。求理论最大连续可达范围，并画出从最早位置到最终位置的一条路径。解释可达与可用的区别。

\begin{solution}
范围为 $1+6(8-1)=43$ 个连续位置，远离序列起点时最早为 $i-42$。一条信息路径依次为输入 $i-42$、第1层 $i-35$、第2层 $i-28$、第3层 $i-21$、第4层 $i-14$、第5层 $i-7$、第6层i。每步最多向后传七格。存在路径不保证信息在有限维状态中不被干扰或压缩损失。
\end{solution}
\end{exercise}

\begin{exercise}\optional 在仅解码器模型的开头加入一个“全局词元”，是否能让它在同一次因果前向中汇总整个后续文档？利用图上的信息方向证明结论。

\begin{solution}
不能。因果图的边从不晚于查询的位置指向该查询，位置0只能读取位置0；由层数归纳，任一层的位置0都不含未来输入。让后续位置读取该全局键只传播初始摘要，不会使它反向汇总全文。可将摘要位置放在片段末端或显式维护外部记忆，但需重新定义协议。
\end{solution}
\end{exercise}

\begin{exercise} 两个长度桶分别为2048和16384，要求实际训练词元占比相等，求按序列抽样的概率。再说明样本重复率为何仍需单独记录。

\begin{solution}
令短序列概率p，要求 $2048p=16384(1-p)$，得 $p=\frac{8}{9}$，长序列概率$\frac{1}{9}$。词元份额相等不代表独立样本量或重复率相等；若某桶可用原始序列很少，有放回抽样会重复使用同一内容，必须记录唯一来源与曝光次数。
\end{solution}
\end{exercise}

\begin{exercise} 使用本章 KV 参数，计算上下文65536、生成上限8192、并发三条时的裸缓存容量，并列出尚未包含的显存项。

\begin{solution}
本章每序列每词元裸KV为128KiB。保留上下文与最大生成总长 $65536+8192=73728$，三条并发需要 $73728\times3\times128\,\mathrm{KiB}=27\,\mathrm{GiB}$。未含模型权重、临时激活和工作区、分页碎片、元数据、调度预留，以及可能的额外推测路径。还要核对模型允许的总长度，不只看输入上限。
\end{solution}
\end{exercise}

\begin{exercise} 设计一个长度、证据位置和干扰程度彼此独立的多证据问题，给出正确答案、无答案变体与评分规则。说明仅测试最长长度的开头位置会遗漏什么。

\begin{solution}
可构造两条证据：“项目甲批准预算70万元”“补充批复增加12万元”，问题为最终批准额，答案82万元；移除任一必须证据形成不可确定变体。独立操纵总长度、两证据相对位置和同格式干扰文件数量，其余文本分布固定。评分同时检查数值、单位、两条来源及无答案拒答。只测最长输入开头会漏掉中部遗漏、跨远距离合并和干扰鲁棒性。
\end{solution}
\end{exercise}

\begin{exercise} 对需要逐条核对长合同数字的任务，比较摘要压缩、检索和直接长窗口三种方法。指出每种方法可能丢失的信息，以及应保留的来源与中间状态。

\begin{solution}
摘要便于总览，但可能删除金额、例外和单位；应保留逐条来源及摘要到原文映射。检索节约上下文，却可能漏掉跨条款定义或补充协议，需版本一致的检索与关联展开。直接长窗口保留更多原文，但仍有位置遗漏和预算限制。可用检索定位加原文精确抽取、单位校验与交叉核对，保留原始条款、查询、候选、计算中间值和不确定项；不可仅依赖自由摘要结算。
\end{solution}
\end{exercise}

\end{exercises}
